\subsection{Giới thiệu}

Hệ thống backend đóng vai trò là trung tâm xử lý chính của ứng dụng \textit{Car Recommendation System}, chịu trách nhiệm tiếp nhận yêu cầu từ người dùng, xử lý nghiệp vụ, truy xuất dữ liệu và trả về kết quả gợi ý phù hợp. Backend hoạt động như một lớp trung gian, kết nối giữa giao diện người dùng (frontend) và các hệ thống lưu trữ dữ liệu phía sau, đảm bảo tính nhất quán, hiệu năng và khả năng mở rộng của toàn bộ hệ thống.

Backend được xây dựng dựa trên framework \textbf{FastAPI} với ngôn ngữ \textbf{Python}, cho phép phát triển các API theo kiến trúc RESTful với hiệu năng cao, khả năng mở rộng tốt và hỗ trợ tài liệu API tự động thông qua OpenAPI/Swagger. FastAPI được lựa chọn nhờ khả năng xử lý bất đồng bộ (asynchronous), phù hợp với các hệ thống cần phục vụ nhiều request đồng thời.

Hệ thống sử dụng \textbf{PostgreSQL} làm cơ sở dữ liệu quan hệ để lưu trữ các dữ liệu có cấu trúc như thông tin xe, người bán, giao dịch và metadata liên quan. PostgreSQL đảm bảo tính toàn vẹn dữ liệu, hỗ trợ các ràng buộc quan hệ và khả năng truy vấn phức tạp. Bên cạnh đó, \textbf{Qdrant} được sử dụng như một vector database, phục vụ cho bài toán tìm kiếm và gợi ý dựa trên vector embedding, cho phép hệ thống thực hiện các phép so sánh tương đồng (similarity search) hiệu quả trên không gian nhiều chiều.

Hệ thống được thiết kế theo hướng \textbf{tách biệt dịch vụ}, trong đó mỗi thành phần đảm nhiệm một vai trò riêng và có thể triển khai, mở rộng độc lập. Trên môi trường sản xuất (Google Cloud Platform), Frontend và Backend là hai dịch vụ \textbf{Cloud Run} riêng biệt: Frontend (React, phục vụ qua Nginx) hiển thị giao diện tại tên miền \texttt{carsalesfinder.com} và ủy quyền các request \texttt{/api} sang Backend; Backend (FastAPI) xử lý nghiệp vụ, kết nối \textbf{AlloyDB} (PostgreSQL 17) qua kết nối SSL và \textbf{Qdrant Cloud} cho tìm kiếm vector. Cloud Run tự động co giãn số instance theo tải, giúp hệ thống không phải tự quản trị hạ tầng máy chủ.

Toàn bộ các service được đóng gói bằng \textbf{Docker}, nhờ đó môi trường phát triển cục bộ có thể tái lập đầy đủ hệ thống bằng Docker Compose (frontend cổng 3000, backend cổng 8000, PostgreSQL cổng 5432, Qdrant cổng 6333) với cùng một image như sản xuất. Cách tiếp cận này đảm bảo tính nhất quán giữa môi trường phát triển và triển khai, đồng thời tạo tiền đề thuận lợi cho việc mở rộng hoặc tích hợp thêm dịch vụ mới trong tương lai.

\subsection{Kiến trúc tổng quan}

\begin{figure}[!htbp]
    \centering
    \includegraphics[width=1\linewidth]{system_architecture_overview.png}
    \caption{Kiến trúc tổng quan hệ thống}
    \label{fig:system-architecture}
\end{figure}

Như minh họa trong Hình \ref{fig:system-architecture}, hệ thống được thiết kế theo kiến trúc phân tầng (layered architecture) kết hợp với mô hình microservices nhằm đảm bảo tính tách biệt chức năng, khả năng mở rộng và dễ bảo trì. Toàn bộ hệ thống được chia thành ba lớp chính: \textit{Client Layer}, \textit{Backend Layer} và \textit{Data Layer}.

\textbf{Client Layer} bao gồm ứng dụng frontend được xây dựng bằng React, được đóng gói tĩnh và phục vụ qua Nginx trên một dịch vụ Cloud Run riêng. Lớp này chịu trách nhiệm hiển thị giao diện người dùng, thu thập dữ liệu đầu vào và gửi các yêu cầu đến backend thông qua các REST API. Client không xử lý logic nghiệp vụ phức tạp mà tập trung vào trải nghiệm người dùng và tương tác.

\textbf{Backend Layer} đóng vai trò trung tâm điều phối và xử lý nghiệp vụ của hệ thống. Backend được xây dựng trên nền tảng FastAPI, triển khai như một dịch vụ Cloud Run độc lập, cung cấp các endpoint RESTful cho frontend. Bên trong backend, các chức năng được chia thành nhiều service độc lập nhằm tuân thủ nguyên tắc single responsibility:
\begin{itemize}
    \item \textit{Authentication Service}: quản lý xác thực người dùng, đăng nhập, phân quyền và kiểm soát truy cập.
    \item \textit{Search Service}: xử lý các yêu cầu tìm kiếm xe dựa trên tiêu chí lọc, từ khóa và các điều kiện truy vấn.
    \item \textit{Recommendation Engine}: thực hiện logic gợi ý xe dựa trên hành vi người dùng và độ tương đồng dữ liệu, kết hợp truy vấn vector embedding.
    \item \textit{Chatbot Service}: cung cấp khả năng tư vấn tự động, hỗ trợ người dùng thông qua hội thoại và tích hợp với mô hình ngôn ngữ lớn.
\end{itemize}

Các service trong backend có thể giao tiếp với nhau thông qua các lời gọi nội bộ và cùng truy cập các nguồn dữ liệu phía dưới, giúp hệ thống dễ dàng mở rộng hoặc bổ sung thêm service mới trong tương lai.

\textbf{Data Layer} chịu trách nhiệm lưu trữ và truy xuất dữ liệu, bao gồm nhiều thành phần chuyên biệt cho từng loại dữ liệu. PostgreSQL (AlloyDB trên môi trường sản xuất) được sử dụng làm cơ sở dữ liệu quan hệ để lưu trữ dữ liệu có cấu trúc như thông tin xe, người dùng và người bán. Qdrant đóng vai trò là vector database, phục vụ cho các bài toán tìm kiếm và gợi ý dựa trên độ tương đồng của vector embedding.

Ngoài ra, hệ thống còn tích hợp với các dịch vụ bên ngoài như OpenAI API để hỗ trợ chatbot và tạo embedding phục vụ cho Recommendation Engine. Việc tách biệt các dịch vụ bên ngoài giúp hệ thống linh hoạt trong việc thay thế hoặc mở rộng các nhà cung cấp khác trong tương lai.

Tất cả các thành phần trong hệ thống được đóng gói dưới dạng container Docker và triển khai trên hạ tầng serverless của Cloud Run, giao tiếp với nhau qua HTTPS. Cách tiếp cận này giúp đảm bảo tính nhất quán giữa các môi trường triển khai, đồng thời tạo nền tảng cho việc mở rộng hệ thống lên các hạ tầng orchestration phức tạp hơn như Kubernetes khi cần thiết.

\subsection{Kiến trúc API}

Backend API được tổ chức theo mô hình phân tầng với bốn lớp chính.

\begin{figure}[!htbp]
    \centering
    \includegraphics[width=1\linewidth]{api_architecture.png}
    \caption{Kiến trúc phân tầng của Backend API}
    \label{fig:api-architecture}
\end{figure}

\noindent \textbf{Middleware Layer} là lớp đầu tiên tiếp nhận request, thực hiện kiểm tra CORS để cho phép cross-origin requests từ frontend, ghi log mọi request để theo dõi và debug, và áp dụng rate limiting để ngăn chặn lạm dụng API.

\noindent \textbf{Authentication Layer} xác thực JWT token trong header Authorization, giải mã và kiểm tra tính hợp lệ của token, sau đó gắn thông tin user vào context để các lớp sau sử dụng.

\noindent \textbf{Router Layer} định tuyến request đến endpoint tương ứng bao gồm auth cho đăng nhập và đăng ký, search cho tìm kiếm xe, reco cho gợi ý, chat cho chatbot, và reviews cho đánh giá. Mỗi router thực hiện validation dữ liệu đầu vào theo schema định nghĩa sẵn.

\noindent \textbf{Service Layer} chứa business logic của từng tính năng. AuthService xử lý hash password và tạo token. SearchService xây dựng query động và cache kết quả. RecommendationEngine tính toán similarity và rank gợi ý. ChatbotService thực hiện RAG workflow với vector database.

\subsection{Thiết kế Database}

Hệ thống sử dụng kiến trúc multi-schema theo mô hình \textbf{Medallion Architecture} với ba tầng dữ liệu: Bronze, Silver và Gold. Các bảng trong Silver và Gold được tạo và quản lý bởi \textbf{dbt} (data build tool), ngoại trừ các bảng user-domain do ứng dụng FastAPI ghi trực tiếp.

\begin{figure}[!htbp]
    \centering
    \includegraphics[width=0.68\linewidth,height=0.68\textheight,keepaspectratio]{database_schema_architecture.png}
    \caption{Kiến trúc multi-schema database}
    \label{fig:database-schema}
\end{figure}

\noindent \textbf{Bronze Schema} lưu trữ dữ liệu thô trực tiếp từ quá trình crawling, chỉ gồm một bảng duy nhất \texttt{bronze.raw\_listings}. Mỗi hàng chứa toàn bộ file JSON crawled dưới dạng cột \texttt{payload} (JSONB), kèm metadata như \texttt{file\_hash} (SHA-256, khóa idempotent), \texttt{gcs\_path}, \texttt{vin}, \texttt{crawl\_date}. Bảng này là \textbf{append-only}: cùng một VIN crawl lại sẽ tạo hàng mới (khác \texttt{file\_hash}), không ghi đè. Việc chọn bản ghi mới nhất được giải quyết ở tầng staging.

\noindent \textbf{Silver Schema} do dbt tạo, tổ chức theo mô hình chiều (Dimensional Model) với các bảng fact/dimension/bridge:
\begin{itemize}
    \item \textbf{Bảng Fact:} \texttt{fct\_listing} (bài đăng xe, incremental delete+insert), \texttt{fct\_model\_rating} (điểm đánh giá tổng hợp), \texttt{fct\_model\_review} (nhận xét chi tiết).
    \item \textbf{Bảng Dimension:} \texttt{dim\_car\_model} (hãng, dòng xe), \texttt{dim\_seller} (thông tin đại lý), \texttt{dim\_feature} (danh mục tính năng), \texttt{dim\_listing\_image} (hình ảnh xe).
    \item \textbf{Bảng Bridge:} \texttt{bridge\_listing\_feature} (quan hệ nhiều-nhiều giữa bài đăng và tính năng).
\end{itemize}

\noindent \textbf{Gold Schema} gồm hai nhóm bảng:
\begin{itemize}
    \item \textbf{Bảng do dbt tạo (marts):} \texttt{vehicles} (denormalized, merge by VIN — bảng chính phục vụ frontend), \texttt{vehicle\_price\_history} (partitioned by month), \texttt{car\_models}, \texttt{sellers}, \texttt{reviews}, \texttt{vehicle\_features}, \texttt{vehicle\_images}. Kèm theo các materialized view: \texttt{mv\_popular\_vehicles} và \texttt{mv\_trending\_models}.
    \item \textbf{Bảng do ứng dụng tạo (user-domain):} \texttt{users}, \texttt{user\_interactions} (hành vi người dùng phục vụ recommendation), \texttt{user\_favorites}, \texttt{user\_searches}, \texttt{user\_reviews}, \texttt{chat\_sessions}, \texttt{chat\_messages}, \texttt{item\_similarity} (ma trận CF precomputed).
\end{itemize}

\textbf{Lưu ý thiết kế:} Trường \texttt{vehicle\_id} trong các bảng user-domain là kiểu TEXT (VIN), \textbf{không có foreign key} đến \texttt{gold.vehicles}. Lý do: bảng \texttt{gold.vehicles} được dbt DROP/CREATE khi chạy full-refresh; một FK cross-schema sẽ phá vỡ quá trình rebuild. Tính toàn vẹn tham chiếu được đảm bảo thay thế bằng dbt test \texttt{assert\_no\_orphan\_interactions}.

\subsection{Data Pipeline}

Quy trình chuyển đổi dữ liệu từ nguồn thô (crawled JSON trên GCS) đến trạng thái sẵn sàng phục vụ ứng dụng được điều phối hoàn toàn tự động bởi \textbf{Temporal} thông qua \texttt{WeeklyPipelineWorkflow}. Chi tiết từng giai đoạn đã được trình bày tại mục \textit{Thu thập và xử lý dữ liệu} ở đầu chương này; phần này tóm tắt luồng xử lý từ góc nhìn backend.

Pipeline gồm ba workflow con chạy tuần tự (fail-stop):

\begin{enumerate}
    \item \textbf{WeeklyCrawlWorkflow:} Thu thập dữ liệu từ Cars.com và upload JSON lên Google Cloud Storage theo cấu trúc phân vùng ngày \texttt{dt=YYYY-MM-DD/}.
    \item \textbf{TransformWorkflow:} Nạp JSON vào \texttt{bronze.raw\_listings} (idempotent qua \texttt{file\_hash}), sau đó chạy \texttt{dbt build} để chuyển đổi qua các tầng staging → silver → gold theo Medallion Architecture.
    \item \textbf{MLWorkflow:} Chạy song song ba tác vụ — tính ma trận \texttt{item\_similarity} (cosine similarity), tạo vector embedding cho hệ thống đề xuất (\texttt{car\_chatbot\_vectors}), và tạo chunked embedding cho chatbot RAG (\texttt{car\_vectorize}).
\end{enumerate}

Toàn bộ pipeline chạy tự động hàng tuần, đảm bảo dữ liệu luôn được cập nhật mà không cần can thiệp thủ công. Tính chất \textbf{incremental} và \textbf{idempotent} cho phép chạy lại an toàn: bronze dedup bằng \texttt{file\_hash}, gold merge by VIN, Qdrant upsert bằng \texttt{point\_id = uuid5(vin)}.

\subsection{Hệ thống Authentication}

Hệ thống sử dụng JWT (JSON Web Token) cho xác thực người dùng.

Khi người dùng đăng ký, hệ thống nhận username, email và password từ frontend. Password được hash bằng bcrypt với salt ngẫu nhiên trước khi lưu vào database. Sau đó tạo JWT token chứa user\_id với thời hạn 30 phút và trả về cho frontend.

Khi người dùng đăng nhập, hệ thống query user từ database theo username, verify password bằng cách so sánh hash, nếu đúng thì tạo và trả về JWT token mới.

Với mỗi request cần xác thực, frontend gửi token trong header Authorization. Backend extract token, decode và verify signature với SECRET\_KEY, kiểm tra expiration time, sau đó query user từ database và gắn vào request context.

\subsection{Tối ưu hiệu năng}

Để đảm bảo hiệu năng, hệ thống kết hợp cache trong tiến trình (in-process) cho các tài nguyên tĩnh và tối ưu truy vấn ở tầng cơ sở dữ liệu.

\subsubsection{Cache trong tiến trình}

Các tài nguyên có chi phí khởi tạo cao nhưng ít thay đổi được nạp một lần và lưu trong bộ nhớ tiến trình bằng \texttt{functools.lru\_cache}, thay vì tính lại ở mỗi request:

\begin{itemize}
    \item \textit{Cấu hình hệ thống gợi ý}: file \texttt{reco\_config.yaml} (trọng số, ngưỡng) được phân tích một lần và cache trong suốt vòng đời tiến trình.
    \item \textit{Tài nguyên chatbot}: LLM client và vector store được khởi tạo một lần khi backend khởi động và tái sử dụng cho mọi lượt hội thoại.
\end{itemize}

Cách tiếp cận này phù hợp với cấu hình \texttt{max-instances = 1} của Cloud Run, giúp giảm độ trễ mà không cần một dịch vụ cache ngoài.

\subsubsection{Tối ưu truy vấn Database}

Các chỉ mục (index) được tạo trên các cột thường xuyên sử dụng trong điều kiện lọc và join. Cụ thể, bảng \texttt{gold.user\_interactions} có index trên \texttt{user\_id}, \texttt{vehicle\_id}, \texttt{interaction\_type} và \texttt{created\_at}. Bảng \texttt{bronze.raw\_listings} có GIN index trên cột \texttt{payload} (JSONB) để tăng tốc truy vấn JSON.

Đối với các truy vấn tổng hợp phức tạp, hệ thống sử dụng \textbf{materialized views} — \texttt{mv\_popular\_vehicles} (xe phổ biến theo lượt tương tác) và \texttt{mv\_trending\_models} (dòng xe được quan tâm nhiều nhất) — được refresh tự động mỗi lần pipeline chạy. Cách tiếp cận này giúp giảm chi phí tính toán tại thời điểm truy vấn và đảm bảo hiệu năng ổn định cho các API thường xuyên được gọi.

\newpage
\subsection{Hệ thống Recommendation}

Recommendation Engine được thiết kế theo kiến trúc \textbf{3-stage hybrid pipeline}, xử lý tuần tự qua ba giai đoạn: candidate generation, ranking và re-ranking. Engine được khởi tạo mới cho mỗi request (stateless) — không fit model trong request, mà sử dụng các tín hiệu precomputed (\texttt{item\_similarity}, vector embeddings) đã tính sẵn bởi pipeline ML.

\begin{figure}[!htbp]
    \centering
    \includegraphics[width=0.8\linewidth]{recommendation_flow.png}
    \caption{Luồng hoạt động của Recommendation Engine}
    \label{fig:reco-flow}
\end{figure}

\subsubsection{Stage 1 — Candidate Generation}

Giai đoạn đầu tiên thu thập một tập ứng viên rộng từ \textbf{4 recaller chạy song song}, mỗi recaller khai thác một nguồn tín hiệu khác nhau:

\begin{itemize}
    \item \textbf{CollaborativeRecaller:} Item-based Collaborative Filtering sử dụng bảng \texttt{gold.item\_similarity} đã precomputed. Với mỗi xe mà user đã tương tác (seed), hệ thống tra cứu K xe láng giềng tương tự nhất. Điểm số được tính có trọng số theo loại tương tác và time decay:
    \[
    \text{score}_{\text{neighbor}} = \sum_{\text{seed}} w_{\text{type}} \cdot e^{-\lambda \cdot \text{days}} \cdot \text{similarity}
    \]
    trong đó trọng số tương tác: view=1.0, click=2.0, compare=3.0, save/favorite=4.0, contact/inquiry=8.0; $\lambda = 0.05$ (half-life $\approx$ 14 ngày).

    \item \textbf{ContentRecaller:} Tìm xe tương tự dựa trên đặc tả kỹ thuật (brand, model, fuel\_type, transmission, drivetrain) và khung giá ($\pm$30\% giá xe seed). Sử dụng SQL query trực tiếp trên \texttt{gold.vehicles}, tính điểm tương đồng theo rule-based scoring:
    \[
    \text{score} = 2.0 \cdot [\text{brand}] + 1.5 \cdot [\text{model}] + 0.5 \cdot [\text{fuel}] + 0.5 \cdot [\text{trans}] + 0.3 \cdot [\text{drive}] + 1.0 \cdot [\text{price\_in\_band}]
    \]

    \item \textbf{VectorRecaller:} Tìm kiếm ngữ nghĩa (semantic similarity) trên Qdrant collection \texttt{car\_chatbot\_vectors}. Truy xuất vector embedding đã lưu của xe seed, sau đó tìm các vector láng giềng gần nhất trong không gian 3072 chiều.

    \item \textbf{PopularityRecaller:} Xe phổ biến toàn hệ thống từ materialized view \texttt{mv\_popular\_vehicles}. Ưu tiên cùng brand với seed vehicle. Đây cũng là \textbf{fallback cho cold-start}: người dùng mới hoặc khách (guest) không có lịch sử tương tác sẽ nhận gợi ý hoàn toàn từ recaller này.
\end{itemize}

Các recaller hoạt động theo cơ chế \textbf{fail-soft}: nếu một nguồn dữ liệu không khả dụng (ví dụ: Qdrant tạm ngừng, matview chưa tồn tại), recaller trả về tập rỗng thay vì raise exception, đảm bảo hệ thống vẫn hoạt động với các nguồn còn lại.

\subsubsection{Stage 2 — Ranking}

Sau khi thu thập ứng viên, module \texttt{FeatureAssembler} hợp nhất kết quả từ 4 recaller, bổ sung thuộc tính xe từ \texttt{gold.vehicles} (brand, model, price, rating, crawled\_at), và tính toán thêm 3 derived features:

\begin{itemize}
    \item \textbf{price\_fit:} Mức độ phù hợp với ngân sách user (= giá trung bình các xe đã xem). Giá càng gần budget, điểm càng cao (1.0 = khớp hoàn toàn, giảm tuyến tính khi lệch).
    \item \textbf{model\_rating:} Điểm đánh giá dòng xe từ Cars.com (0--5, normalize về 0--1).
    \item \textbf{recency:} Độ mới của bài đăng, sử dụng exponential decay với half-life $\approx$ 30 ngày.
\end{itemize}

\texttt{WeightedLinearRanker} tính điểm tổng cho mỗi ứng viên bằng công thức:
\[
\text{rank\_score} = \sum_{i=1}^{7} w_i \cdot f_i
\]
với 7 features và trọng số cấu hình trong \texttt{reco\_config.yaml}:

\begin{table}[H]
\centering
\small
\caption{Bảy đặc trưng và trọng số của WeightedLinearRanker}
\label{tab:ranker-features}
\begin{tabular}{|l|c|l|}
\hline
\textbf{Feature} & \textbf{Weight} & \textbf{Nguồn} \\
\hline
cf\_score & 3.0 & CollaborativeRecaller \\
\hline
vector\_score & 2.0 & VectorRecaller (Qdrant) \\
\hline
content\_score & 1.5 & ContentRecaller (SQL) \\
\hline
popularity & 1.0 & PopularityRecaller (matview) \\
\hline
price\_fit & 1.0 & Derived (budget proximity) \\
\hline
model\_rating & 0.8 & Cars.com consumer rating \\
\hline
recency & 0.5 & Listing freshness \\
\hline
\end{tabular}
\end{table}

\textbf{CF warmup:} Trọng số \texttt{cf\_score} được nhân thêm hệ số \texttt{cf\_scale} $\in [0, 1]$ tỉ lệ thuận với số tương tác của user. Khi user mới (0 interactions), $\text{cf\_scale} = 0$ và CF không đóng góp vào điểm. Khi user đạt ngưỡng 20 interactions, $\text{cf\_scale} = 1.0$ và CF phát huy tối đa. Cơ chế này giúp hệ thống chuyển dịch tự nhiên từ content/popularity-driven sang collaborative-driven khi tích lũy đủ dữ liệu hành vi.

Ranker được thiết kế theo mô hình \textbf{weighted-linear} (thay vì black-box model) để đảm bảo tính giải thích được (explainability): có thể chỉ ra chính xác tín hiệu nào đẩy một xe lên cao trong danh sách.

\subsubsection{Stage 3 — Re-ranking}

\texttt{MMRReranker} (Maximal Marginal Relevance) cân bằng giữa \textbf{relevance} và \textbf{diversity} trong kết quả cuối cùng:
\[
\text{MMR} = \lambda \cdot \text{relevance} - (1 - \lambda) \cdot \max_{\text{đã chọn}} \text{similarity}
\]
với $\lambda = 0.7$ (nghiêng về relevance nhưng vẫn phạt trùng lặp). Similarity giữa hai xe được tính nhanh theo brand (0.5) và model (0.5) mà không cần embedding.

Ngoài ra, hệ thống áp dụng \textbf{hard caps} để tránh kết quả đơn điệu:
\begin{itemize}
    \item Tối đa \textbf{3 xe} cùng brand trong danh sách cuối.
    \item Tối đa \textbf{2 xe} cùng car\_model.
\end{itemize}

Khi tất cả ứng viên còn lại đều vi phạm caps, hệ thống \textit{relax} ràng buộc và chọn theo relevance để đảm bảo đủ \texttt{top\_k} kết quả (mặc định 20).

\subsubsection{Tóm tắt luồng xử lý}

\begin{table}[H]
\centering
\small
\caption{Tóm tắt ba giai đoạn của Recommendation Engine}
\label{tab:reco-stages}
\begin{tabular}{|c|l|l|}
\hline
\textbf{Stage} & \textbf{Module} & \textbf{Mô tả} \\
\hline
1 & candidates.py & 4 recaller song song $\rightarrow$ tập ứng viên (pool\_size $\leq$ 300) \\
\hline
2 & features.py + ranker.py & Bổ sung 7 features, tính rank\_score tuyến tính \\
\hline
3 & reranker.py & MMR diversity + brand/model caps $\rightarrow$ top\_k = 20 \\
\hline
\end{tabular}
\end{table}

Toàn bộ trọng số và ngưỡng được cấu hình tập trung trong file \texttt{reco\_config.yaml}, không có magic numbers trong code. Kiến trúc này cho phép điều chỉnh chiến lược gợi ý chỉ bằng cách thay đổi file YAML mà không cần sửa code.

